\section{The ten-integer companion}\label{sec:ten}

This section gives the complete evaluator over $\Z^{10}$. It uses the
\emph{unshifted} transport row in the table of Section~\ref{sec:sparse},
the same support-isolating labels, and the existing package-40/Sun
arithmetic bridge. The two constructions therefore use slightly different
parent coefficient records, despite having the same supported pairs.
The ten unrestricted coordinates are exactly
\[
f,g,h,k,l,w,x,y,z,m.
\]
Every other symbol introduced in this section is a polynomial
abbreviation, not an additional witness.

\subsection{Fixed source and the required range contract}

Let $\Dcode$ be the factorial-weighted kernel of the unshifted parent,
with the targets $E_j=400+401(j-1)$. Thus its degree is 8821 and its
leading coefficient is 2. Keep the whole package-31 tuple and put
\[
 H_0=8(1+\beta+K_0+M_C+M_F)^2,
 \qquad \beta_S=2^{2140030}.
\]
The ordinary quadratic coefficients have magnitude at most $H_0$,
and the factorial-weighted ones have magnitude at most $2H_0$.
For example, their port degree is at most two and their integer
coefficient $\ell^1$ norm is at most eight. The supplied coefficient
audit verifies these statements exactly. The port bit lengths give
\[
 H_0<2^{2140017},\qquad
 \beta_S>4\cdot33^2H_0.
\]
The actual integers and all sixteen live nonliteral coefficient leaves
are materialized in the archive. These coefficient choices are fixed;
they do not add unknowns to the equation.

Define the source registers by
\begin{align*}
 b&=1+3f,&\BB&=\beta_Sb^2,&
 \ell&=\sum_{r=1}^{31}\BB^{w_r},\\
 N_0&=\BB^{201},&N_1&=4\BB^{9222},&N&=N_0N_1=4\BB^{9423},\\
 M_c&=N_0-1-(b-1)\ell,&
 T_*&=\sum_{r=1}^{22}\BB^{E_r},\\
 T_c&=M_c+(\BB-2)T_*N_0,\\
 J_c&=2(1+g)^2\Dcode(\BB)+\BB\sum_{r=0}^{9221}\BB^r,\\
 S_c&=g+J_cN_0,\\
 R_c&=(S_c+T_c+1)N+T_c+1,\\
 X&=(N-1)R_c,&Y&=N^2,&\mathcal C&=b\BB^{200}.
\end{align*}
In particular, the upper-range polynomial has the exact fixed form
\begin{equation}\label{eq:tencap}
 \mathcal C=2^{428006000}b^{401}.
\end{equation}
Its power-of-two exponent is positive and even, and its exponent of
$b$ is positive and odd. All exponents in the source are fixed
nonnegative integers.

\begin{proposition}[Source contract for the ten-integer bridge]
\label{prop:tensource}
The displayed source has the following properties.
\begin{enumerate}[label=\textup{(\roman*)},leftmargin=*]
\item From a common natural zero of the unshifted parent one can choose
 arbitrarily large $f>0$, with $b$ both a square and a power of two,
 and an integer $g$ satisfying
 \[
 b\le g<\mathcal C,\qquad Y\mid\binom{2X}{X}.
 \]
\item If $f\ne0$, $b$ is both a square and a power of two,
 $0\le g<2\mathcal C$, and $Y\mid\binom{2X}{X}$, then the
 parent has a common natural zero.
\item If $f\ne0$, $b$ is a square and $0\le g<2\mathcal C$,
 without yet assuming a power-of-two base, then
 \[
 0\le S_c,T_c<N,\qquad 3\mid X,\qquad
 X>3b,\qquad Y>\max\{b,256\}.
 \]
\end{enumerate}
\end{proposition}
\begin{proof}
Under the hypotheses of (iii), $b=1+3f$ forces $f>0$ and $b\ge4$.
The positive leading coefficient and the bound $2H_0$ on the other
kernel coefficients imply $\Dcode(\BB)>0$. Integrality and
$g<2b\BB^{200}$ give $1+g\le2b\BB^{200}$ and $g<N_0$.
Since $16H_0b^2\le\BB-2$, the kernel estimate yields
\[
\begin{split}
0<J_c
&\le (16H_0b^2+\BB)\sum_{r=0}^{9221}\BB^r\\
&<2\BB^{9222}<N_1.
\end{split}
\]
The mask satisfies $0\le M_c<N_0$, while
$(\BB-2)T_*<\BB^{8822}<N_1$. Hence both packed words lie in
$[0,N)$. The exponent in $\beta_S$ is even, so
$\beta_S\equiv b\equiv\BB\equiv1\pmod3$ and $N\equiv1\pmod3$.
Therefore $3\mid X$. Also $R_c\ge N+1>\BB>b$ and $N\ge4$,
giving $X>3b$; the fixed base makes $Y>\max\{b,256\}$ immediate.

If $b$ is a power of two, the words $\BB,N_0,N$ are powers of two.
Lemma~\ref{lem:carry} identifies the binomial divisibility with the
joint no-carry test of $S_c,T_c$. Its low block forces
$g=\sum_r n_r\BB^{w_r}$ with $0\le n_r<b$. Theorem~\ref{thm:extract}
recovers the twenty-two residuals at the target cells. Every coefficient
of the square times the kernel has absolute value at most
\[
 2H_0(1+\textstyle\sum_r n_r)^2<\BB/2,
\]
so the geometric offset and target masks test precisely whether all
residuals vanish. This proves (ii).

Conversely, take a parent zero and choose $b=2^{2q}$ as large as needed,
with $b$ exceeding all its coordinates. Then $f=(b-1)/3$ is an arbitrarily
large positive integer. Pack the zero into $g$. It is not the all-zero
tuple, and all selected weights are positive, so $g\ge\BB\ge b$.
The geometric bound gives $g<b\BB^{200}=\mathcal C$.
The same masks pass, and the binomial divisibility follows. This proves
(i) and completes the contract.
\end{proof}

\subsection{The complete polynomial evaluator}

Continue with the following substitutions:
\begin{align*}
 L&=lY,&U&=2LX,\\
 V&=4gwY,&W_b&=bw,\\
 A&=U(V+1),&B_1&=2X+1,\\
 C&=B_1+(A-2)h,&P_n&=U^2V,\\
 K&=X+1+k(P_n-2),&S&=2A-5,\\
 T&=3W_bC-2(W_b^2-1),&R_2&=(P_n^2-4)K^2+4.
\end{align*}
\begin{align*}
 \Delta_b&=8\mathcal C^3gK^2
       -g^2\bigl(32(C-KL)^2\mathcal C^3+g^2K^2\bigr),\\
 O&=f^2l^2\Delta_b,\qquad \xi=zx.
\end{align*}
The modified Lucas-certificate registers are
\begin{align*}
 D&=(A^2-4)C^2+4,&E&=C^2D\xi,\\
 F&=4(A^2-4)E^2+1,&G_*&=2(A+2)(A-2)^2E^2-1,\\
 \omega&=16E,&H_*&=C+(B_1-C)F+\omega Fy,\\
 I_*&=(G_*^2-1)H_*^2+1.
\end{align*}
Define the divisibility merger, radicands and separating weights by
\begin{align*}
 \Sigma&=\xi S^5,&r&=T+Sy,\\
 \Theta&=(r^3-rS^2-S^3)(r^2+23S^2),\\
 r_0&=b\bigl(1+(3O-4)b^2z^2\bigr),&
 r_1&=R_2,&r_2&=DFI_*,\\
 v_0&=1,&v_1&=r_0,&v_2&=r_0(P_n^2+K^2+3).
\end{align*}
Finally put
\begin{align*}
 Z_n&=\Sigma m-\Theta,&a_i&=\Sigma^2v_i^2r_i\quad(0\le i\le2),\\
 V_n&=Z_n^2+a_2-a_0-a_1,\\
 W_n&=V_n^2+4Z_n^2a_2-4a_0a_1.
\end{align*}
The requested unrestricted-integer equation is
\begin{equation}\label{eq:F10}
 \boxed{F_{10}=W_n^2-16V_n^2Z_n^2a_2=0.}
\end{equation}
Together with the fixed coefficient records and the source above, this
specifies every arithmetic operation needed to define one polynomial
in the ten displayed integer unknowns. There are no divisibility
predicates, roots, variable exponents or hidden existential parameters
in its definition.

\subsection{The imported parameterized arithmetic bridge}

We state the precise inherited interface rather than silently applying
a theorem whose literal statement uses a different source polynomial.
Package 40 reads the proofs of Sun's Propositions 6.1--6.2 parametrically:
the source enters through Proposition~\ref{prop:tensource},
$b=1+3f$, and a bound of the form $\mathcal C=2^c b^d$ with fixed
positive even $c$ and positive odd $d$. Its size and congruence
hypotheses are exactly those proved in part (iii). The new constants
$c=428006000$ and $d=401$ satisfy that interface. The polynomial
tail and the arithmetic witnesses have not changed
\cite{package40,sun2021,sun2026}.

For clarity, let $u_0(A)=0$, $u_1(A)=1$ and
$u_{q+1}(A)=Au_q(A)-u_{q-1}(A)$, extended in the usual way to
integer indices. The imported bridge and interpolation statements are:
\begin{enumerate}[leftmargin=*]
\item A parent zero admits $f,g,h,k,l,w$ with $A>2$ even,
 $B_1>0$ odd, $b$ a square, $b\le g<\mathcal C$,
 $C=u_{B_1}(A)$, $R_2$ a square, $S\mid T$, and $O>0$.
\item For the displayed substitutions, the conditions
 \[
 b\text{ square},\quad O>0,\quad \xi\ne0,\quad
 DFI_*\text{ square},\quad R_2\text{ square},\quad S\mid T
 \]
 imply a parent zero.
\item If $A>2$ is even, $B_1>0$ is odd and $C=u_{B_1}(A)$,
 then for every positive integer $d_0$ there is $\xi>0$ divisible
 by $d_0$ such that each residue class modulo $\xi$ contains a $y$
 for which $DFI_*$ is a square.
\end{enumerate}
The third statement is Sun's full residue-class interpolation theorem.
It is needed because the divisibility merger below prescribes a residue
class after choosing $\xi$. We retain this theorem as a cited dependency;
we do not claim a fresh proof of the Lucas results or of the general
ten-integer theorem.

The reverse direction does not assume a power-of-two base prematurely.
Once $b$ is a square and $O>0$, the displayed formula for $O$ forces
$f,l\ne0$, $\Delta_b>0$, $g>0$ and $K\ne0$. It then gives
\[
 0<g<2\mathcal C,\qquad 4(C-KL)^2<K^2.
\]
Part (iii) of Proposition~\ref{prop:tensource} is therefore available
before the Lucas step. The inherited reverse argument recovers
\[
 C=u_{B_1}(A),\quad K=u_{X+1}(P_n),\quad
 W_b=2^{B_1},\quad Y\mid\binom{2X}{X}.
\]
Since $W_b=bw$ and $b>0$, it follows that $b$ is a power of two.
Only at this point is part (ii) applied. This order of reasoning is
why the broad pre-decoding range contract was proved separately.

\subsection{Why the finalizer introduces no extra solutions}

Here we recall the short algebraic part of the inherited finalizer,
including its nonvanishing and choice-order conditions. The integer
$S=2A-5$ is always odd and nonzero. The polynomial
\[
 \pi(v)=(v^3-v-1)(v^2+23)
\]
has no rational root and has a root modulo every positive integer.
For the latter assertion, at an odd prime other than 23 either the
cubic has a simple root or it is irreducible; in the irreducible case
Frobenius is a three-cycle, so its discriminant $-23$ is a square and
the quadratic has a simple root. At 23 the cubic has the simple root
3. These roots lift to odd prime powers. At 2, the fact
$-23\equiv1\pmod8$ supplies a root in the 2-adic integers.
The Chinese remainder theorem finishes the claim. The rational-root
test for the cubic and positivity of the quadratic prove nonvanishing
over $\Q$.

If $S\mid T$, write $T=Sq$. For any $\xi\ne0$ choose a root
$c_0$ of $\pi$ modulo $|\xi|$ and prescribe
$y\equiv c_0-q\pmod\xi$. Then $\Sigma\mid\Theta$.
Conversely, $S^5\mid\Theta$ forces $S\mid T$: if
$\nu_p(S)>\nu_p(T)$, the two homogeneous factors of $\Theta$
have valuations $3\nu_p(T)$ and $2\nu_p(T)$, contradicting
$S^5\mid\Theta$. Also $\Theta\ne0$, since otherwise
$\pi((T+Sy)/S)=0$.

Suppose $z\ne0$ and $r_0$ is a square. The factors
$b$ and $1+(3O-4)b^2z^2$ are coprime. The second cannot be zero:
that would force $|b|=|z|=1$, hence $b=1$, $f=O=0$, contradicting
its vanishing. Unique factorization makes $b$ a signed square;
$b\equiv1\pmod3$ excludes the negative sign. The other factor is
then a square as well, which forces $O>0$ because for $O\le0$
it is at most $1-4(bz)^2<0$.

For the converse, the positive bridge gives $O\ge g\ge b\ge4$:
$\Delta_b$ is a positive multiple of $g$, and $f,l$ are nonzero.
Thus $d=3O-4$ is positive and is 2 modulo 3, so it is nonsquare.
A positive Pell unit for $d$, raised to a suitable finite order
modulo $b$, supplies
$u^2-dv^2=1$ with $v>0$ and $b\mid v$.
Choose $z=v/b>0$. Then $r_0$ is a square. The interpolation theorem
can now choose $\xi>0$ divisible by this already fixed $z$; put
$x=\xi/z$. Only after this modulus exists do we prescribe the
residue class of $y$ needed for $\Sigma\mid\Theta$.

For every integer assignment with $z\ne0$, one has $|r_0|>1$.
Indeed $r_0=0,\pm1$ would force $|b|=1$, hence $b=1$ and
$f=O=0$, whereas $r_0=1-4z^2\le-3$.
Also
\[
 \sqrt{|r_1|}\le P_n^2+K^2+2,
\]
as follows by subtracting $(P_n^2+4)K^2+4$ from
$(P_n^2+K^2+2)^2$. Consequently
\[
 |v_1|>\sqrt{|r_0|},\qquad
 |v_2|>\sqrt{|r_0|}+|v_1|\sqrt{|r_1|}.
\]
If a signed sum of these three weighted roots is rational, every
radicand is an integer square. To see this, in the multiquadratic
extension take the largest-index nonzero root term changed by an
automorphism. Its change has greater absolute value than all preceding
changes combined, a contradiction. This argument works with complex
absolute values for negative radicands and does not assume independent
square classes. Rational square roots of integers are integral.

The final polynomial has the exact identity
\[
 F_{10}=\prod_{\varepsilon\in\{\pm1\}^3}
 \left(Z_n+\varepsilon_0\sqrt{a_0}
           +\varepsilon_1\sqrt{a_1}
           +\varepsilon_2\sqrt{a_2}\right).
\]
The radicals cancel; successive conjugate multiplication gives precisely
\eqref{eq:F10}. If $\Sigma=0$, then all $a_i=0$ and
$F_{10}=\Theta^8\ne0$. At any zero, therefore, $\xi=zx\ne0$.
Divide a vanishing factor by $\Sigma$. Separation makes all $r_i$
squares and the resulting signed sum integral, so $\Sigma\mid\Theta$.
The two mergers give $S\mid T$, $b$ square and $O>0$; the reverse
bridge then yields a parent zero.

Conversely, a parent zero gives the positive bridge witnesses. Choose
$z$, $\xi$, $x$ and $y$ in the order just described, obtaining all
three square radicands and $\Sigma\mid\Theta$. An integer $m$ can
then make $Z_n$ equal to the negative of $\Sigma$ times any signed
sum of the weighted integer roots. One norm factor vanishes.
This proves both implications for the ten-integer equation and hence
\eqref{eq:logicaltarget}, using the explicitly stated inherited
compiler and arithmetic contracts.

\subsection{Exact degree and evaluator size}

Put $u=9221+200+2=9423$. The geometric offset uniquely dominates
$J_c$, and the source registers have degrees
\[
\begin{array}{c|rrrrrr}
\text{register}&N&S_c&T_c&X&Y&\mathcal C\\\hline
\text{degree}&2u&2u&2u-800&6u&4u&401.
\end{array}
\]
In particular $d_X=56538$, $d_Y=37692$ and $d_C=401$.
Propagation through the retained tail gives
\begin{align*}
 \deg Z_n&=5d_X+10d_Y+30=659640,\\
 \deg a_0&=14d_X+28d_Y+3d_C+57=1848168,\\
 \deg a_1&=26d_X+48d_Y+6d_C+98=3281708,\\
 \deg a_2&=114d_X+224d_Y+6d_C+440=14891186.
\end{align*}
The last degree exceeds those of $a_0,a_1,Z_n^2$. Therefore the
highest part of $F_{10}$ is the fourth power of the highest part
of $a_2$, giving
\[
 \deg F_{10}
 =456d_X+896d_Y+24d_C+1760
 =59\,564\,744.
\]
The independent actual-coefficient certificate also gives
\[
 [T^{59564744}]F_{10}(T,\ldots,T)
 \equiv285435419\pmod{1000000007},
\]
which is nonzero. This is an exact leading-coefficient computation,
not sampling of possible solutions.

The primary division-free graph has 215 additions, 25 subtractions
and 413 multiplications, totaling \textbf{653 operations}.
Shared powers, row-wise quadratic Horner evaluation and divide-and-
conquer geometric sums account for its size. Every binary gate is
charged under the common convention. These are exact construction
metrics and do not assert that no shorter ten-integer evaluator exists.
